%Łukasz DZIEŁ
%tel: +48 883 533 374
%email: lukasz.dziel@wat.edu.pl
%wersja: 20220316
%testowano w środowisku: TeXLive2021 (instalacja pełna z iso) + TeXstudio 4.2.2

\documentclass[13pt, a4paper, twoside]{scrartcl}

\usepackage[left=20mm,right=20mm,top=25mm,bottom=25mm,bindingoffset=10mm]{geometry}
\usepackage[T1]{fontenc}
\usepackage[polish]{babel}
\usepackage[utf8]{inputenc}

\usepackage{mathptmx}
\usepackage{courier}

\usepackage[normalem]{ulem}

\usepackage{sectsty}
\allsectionsfont{\normalfont\bfseries\fontsize{14}{18}\selectfont}

\usepackage[titles]{tocloft}
\setcounter{tocdepth}{2}
\renewcommand{\cftsecfont}{\fontsize{12}{14}\selectfont}
\renewcommand{\cftsecpagefont}{\fontsize{12}{14}\selectfont}
\renewcommand{\cftsecleader}{\fontsize{12}{14}\selectfont\cftdotfill{1}}
\renewcommand{\cftsecindent}{0mm}
\renewcommand{\cftsecnumwidth}{22mm}
\renewcommand{\cftsecaftersnum}{. }
\renewcommand{\cftbeforesecskip}{2mm}

\renewcommand{\cftsubsecfont}{\fontsize{11}{13}\selectfont}
\renewcommand{\cftsubsecpagefont}{\fontsize{11}{13}\selectfont}
\renewcommand{\cftsubsecindent}{5mm}
\renewcommand{\cftsubsecnumwidth}{8mm}
\renewcommand{\cftsubsecaftersnum}{. }
\renewcommand{\cftbeforesubsecskip}{2mm}
\renewcommand{\cftsubsecleader}{\fontsize{11}{13}\selectfont\cftdotfill{1}}

\renewcommand{\cftsubsubsecfont}{\fontsize{11}{13}\selectfont}
\renewcommand{\cftsubsubsecpagefont}{\fontsize{11}{13}\selectfont}
\renewcommand{\cftsubsubsecindent}{5mm}
\renewcommand{\cftsubsubsecnumwidth}{12mm}
\renewcommand{\cftsubsubsecaftersnum}{. }
\renewcommand{\cftbeforesubsubsecskip}{0mm}
\renewcommand{\cftsubsubsecleader}{\fontsize{11}{13}\selectfont\cftdotfill{1}}

\renewcommand{\cftfigindent}{0mm}
\renewcommand{\cftfignumwidth}{15mm}
\renewcommand{\cftfigpresnum}{Rys. }
\renewcommand{\cftfigaftersnum}{. }
\renewcommand{\cftfigleader}{\normalfont\cftdotfill{1}}

\renewcommand{\cfttabindent}{0mm}
\renewcommand{\cfttabnumwidth}{15mm}
\renewcommand{\cfttabpresnum}{Tab. }
\renewcommand{\cfttabaftersnum}{. }
\renewcommand{\cfttabfont}{\normalfont}
\renewcommand{\cfttableader}{\normalfont\cftdotfill{1}}

\usepackage{amsmath,amsfonts,amsthm}

\usepackage{color}
\usepackage{float}

\usepackage{verbatimbox}
\def\verbarg{{\fontsize{11}{13}\selectfont\makebox[7mm]{\arabic{VerbboxLineNo}}}\hspace{12mm}}

\usepackage[plain]{algorithm}
\floatstyle{plaintop}
\restylefloat{algorithm}
\floatname{algorithm}{Alg.}

\usepackage[pdftex]{graphicx}

\renewcommand\thesection{Rozdział \Roman{section}}
\renewcommand\thesubsection{\Roman{section}.\arabic{subsection}}
% \renewcommand\thesubsubsection{\Roman{section}.\arabic{subsection}.\arabic{subsubsection}}

\usepackage{longtable}
\usepackage[font=bf, labelfont=bf, belowskip=-0mm]{caption}
\captionsetup[figure]{name=Rys., labelsep=period}
\captionsetup[table]{name=Tab.,  labelsep=period}
\captionsetup[lstlisting]{margin=0pt, font=bf, labelsep=period}
\captionsetup[algorithm]{font=bf, labelsep=period}

\usepackage{listings}
\usepackage{color}

\usepackage{url}
\urlstyle{ttfamily}


\newcommand{\myparagraph}[1]{\paragraph{#1}\mbox{}\\}
\newcommand{\myurl}[1]{\urlstyle{same}\url{#1}}

\usepackage{hyperref}


\usepackage{fancyhdr}
\pagestyle{fancy}
\fancyhead[RO, LE]{\thepage}
\fancyhead[RE, LO]{}
\fancyfoot{}
\renewcommand\headrulewidth{0pt}

\fancypagestyle{firststyle}{
\fancyhead[RO, LE]{}
\fancyhead[RE, LO]{}
}

\setlength\parindent{27pt}

\newcommand{\inserttitlepage}{
    \thispagestyle{firststyle}
    \begin{center}
        \fontsize{22}{14}\textbf{WOJSKOWA~~AKADEMIA~~TECHNICZNA}\\
        \fontsize{12}{10}\textbf{im. Jarosława Dąbrowskiego}
    \end{center}
    \vspace{-7mm}
    \rule{\linewidth}{0.5pt}
    \begin{center}
        \fontsize{20}{10}{\textbf{WYDZIAŁ~~CYBERNETYKI}}
    \end{center}
    
    \vspace{5mm}
    
    \begin{center}
        \includegraphics[width=40mm]{./defs/WAT.png}
    \end{center}

    \vspace{0mm}

    \begin{center}
        {{\fontsize{26}{10} \selectfont {PRACA DYPLOMOWA}}}\\
        \vspace{3mm}
        {\fontsize{20}{10} \selectfont \stopien}
    \end{center}
    
    \vspace{5mm}
    
    \begin{longtable*}{p{.18\textwidth} p{.75\textwidth}}
        \fontsize{14}{20}\selectfont{Temat pracy:} &\fontsize{16}{20}\selectfont{\textbf{\noindent\temat}} \\
    \end{longtable*}
    
    \vspace{5mm}
    
    \begin{center}
        \textrm{\textbf{\fontsize{14}{20}\selectfont{\kierunek}}}\\
        \scriptsize\textrm{...........................................................................................................\\{{\scriptsize{(kierunek studiów)}}}}\\~\\
        \vspace{5mm}    
        \textrm{\textbf{\fontsize{14}{20}\selectfont{\specjalnosc}}}\\
        \scriptsize\textrm{...........................................................................................................\\\scriptsize{(specjalność)}}\\
    \end{center}
   
    \vspace{5mm}

    \begin{longtable*}{p{.5\textwidth} p{.5\textwidth}}
        \textrm{Dyplomant:}\vspace{5mm} & 	\textrm{Promotor:}\\
        \textrm{\textbf{\autor}}  & \textrm{\textbf{\promotor}} \\
    \end{longtable*}
    
    \begin{center} 
    \rule{\linewidth}{0.5pt}
    \fontsize{12}{10} \textbf{\data} \end{center}
    
     
    \clearpage
    
    \thispagestyle{firststyle}
    
    \fontsize{14}{20} \selectfont{\textbf{OŚWIADCZENIE}}\\
    
    \begin{center}
        \fontsize{14}{20} \selectfont \textit{\ifdefined\zgoda  Wyrażam zgodę / \sout{nie wyrażam zgody} \else \sout{Wyrażam zgodę} / {nie wyrażam zgody} \fi $^{*}$\\na udostępnianie mojej pracy przez Archiwum WAT}
    \end{center}
    
    ~\\
    
    \begin{longtable*}{p{.6\textwidth} c}
    Dnia ...................... & 	....................................\\
    ~ & \fontsize{10}{10} \selectfont (podpis)\\
    \end{longtable*}
    \textit{$^{*}$ Niepotrzebne skreślić}
    \newpage
    \normalsize
    
    \tableofcontents
    
    \newpage
}

\defcaptionname{polish}{\refname}{Bibliografia}

\definecolor{mygray}{rgb}{0.95,0.95,0.95}

\renewcommand{\lstlistingname}{Kod.}

\lstset{
    basicstyle=\ttfamily\footnotesize,
    numbers=left,
    numberstyle=\ttfamily\footnotesize\color{black},
    numbersep=5pt,
    showspaces=false,
    breaklines=true,
    keepspaces=true,
    tabsize=2,
    xleftmargin=20pt,
    framexleftmargin=20pt,
    captionpos=tc,
    backgroundcolor=\color{mygray},
    keywordstyle=\color{black},
    inputencoding=utf8
}

\newcommand{\source}[1]{\vspace{0mm} \noindent\fontsize{10}{10}\selectfont{Źródło: {#1}}\normalsize\vspace{0mm}}

\newcommand{\TODO}[1]{
    {\large \color{red}{\underline{\textbf{TODO: #1}}}} 
}